(1) 基频诱导的二阶极化随空间相位 \(2k(\omega)z\) 变化，而自由传播的倍频模相位为 \(k(2\omega)z\)，差值 \(\Delta k=2k(\omega)-k(2\omega)\) 导致沿晶体不同位置产生的倍频光无法持续同相累加。题给慢变包络方程已经固定复幅的相位约定与耦合常数 \(\kappa\)。
(2) 积分 \(A_2(L)=i\kappa A_1^2\int_0^L e^{i\Delta kz}dz\)。用指数积分或中点相位展开，得 \(A_2(L)=i\kappa A_1^2L\,e^{i\Delta kL/2}\operatorname{sinc}(\Delta kL/2)\)，其中 \(\operatorname{sinc}x=\sin x/x\)，零处以连续极限定义。
(3) 倍频输出强度在相同光阻抗约定下与 \(|A_2|^2\) 成正比，因此 \(I_{2\omega}(L)\propto|\kappa|^2|A_1|^4 L^2\operatorname{sinc}^2(\Delta kL/2)\)。若 \(\Delta k=0\)，\(\operatorname{sinc}(0)=1\)，得到 \(I_{2\omega}\propto L^2\)，表明场幅随长度线性相干叠加。
(4) 相位差 \(\Delta kz\) 首次达到 \(\pi\) 的长度为 \(L_c=\pi/|\Delta k|\)；倍频振幅的第一非平凡零点要求 \(\Delta kL/2=\pi\)，故 \(L_{\rm zero}=2\pi/|\Delta k|=2L_c\)。当倍频转换已显著消耗基频能量时，\(A_1\) 不再常量，必须联立耦合方程并计入能量转移。
